\documentclass[11pt]{article}
\usepackage[margin=1in]{geometry}
\usepackage{enumitem}
% =============================================================================
% Preambles/header.tex — shared LaTeX/Beamer preamble for this project
%
% Use from a Beamer lecture by prepending:
%     \input{header}
% and compiling with TEXINPUTS=../Preambles:$TEXINPUTS (the /compile-latex
% skill does this automatically).
%
% PALETTE CONTRACT. Color names defined here MUST match the names in
% Quarto/theme-template.scss so Beamer and Quarto renderings use the same
% palette. The scripts/check-palette-sync.sh script enforces this.
%
% To customize: edit the HEX values below AND the matching $variable in
% Quarto/theme-template.scss, then run scripts/check-palette-sync.sh.
% =============================================================================

% -----------------------------------------------------------------------------
% Engine detection + encoding
%
% Our /compile-latex skill uses XeLaTeX, where inputenc is unnecessary and
% can produce encoding warnings. Only load inputenc under pdfTeX.
% -----------------------------------------------------------------------------
\usepackage{iftex}
\ifPDFTeX
  \usepackage[utf8]{inputenc}
\fi

\usepackage{amsmath, amssymb, amsthm}
\usepackage{booktabs}
\usepackage{graphicx}

% -----------------------------------------------------------------------------
% PALETTE — mirrors Quarto/theme-template.scss variable names.
% Load xcolor and define colors BEFORE anything that references them
% (notably hyperref, hypersetup, and the Beamer color assignments).
% -----------------------------------------------------------------------------
\usepackage{xcolor}

\definecolor{primary-blue}{HTML}{012169}      % headings, accents
\definecolor{primary-gold}{HTML}{B9975B}      % emphasis, borders
\definecolor{highlight-yellow}{HTML}{F2A900}  % markers, alerts
\definecolor{light-bg}{HTML}{E8EDF5}          % title / section backgrounds
\definecolor{jet}{HTML}{1A1A1A}               % body text

% Semantic colors (used in diagrams and annotations)
\definecolor{positive}{HTML}{15803D}          % good / observed / identified
\definecolor{negative}{HTML}{B91C1C}          % bad / problematic / bias
\definecolor{neutral}{HTML}{525252}           % reference / context

% Highlight accents (match .hi-* classes in the SCSS)
\definecolor{hi-slate}{HTML}{314F4F}
\definecolor{hi-green}{HTML}{2E8B57}
\definecolor{hi-red}{HTML}{C41E3A}

% Semantic aliases so skill templates and snippets can write
% \textcolor{accent}{...} without hard-coding "primary-blue".
\colorlet{accent}{primary-blue}
\colorlet{accent2}{primary-gold}

% -----------------------------------------------------------------------------
% Hyperref — loaded AFTER xcolor + palette so link colors resolve.
% -----------------------------------------------------------------------------
\usepackage{hyperref}
\hypersetup{colorlinks=true, linkcolor=primary-blue, urlcolor=primary-blue,
            citecolor=primary-blue, pdfborder={0 0 0}}

% -----------------------------------------------------------------------------
% Course code — set once per deck (after \input{header}, before \begin{document})
% via \coursecode{CS401}. Shown in the footer on every slide so a deck's course
% is identifiable without opening the title page. Empty by default.
% -----------------------------------------------------------------------------
\makeatletter
\newcommand{\coursecode}[1]{\gdef\@coursecodetext{#1}}
\gdef\@coursecodetext{}
\makeatother

% -----------------------------------------------------------------------------
% Beamer theme assignments (only apply under Beamer).
%
% We detect Beamer via \@ifclassloaded{beamer}, which is the documented,
% non-internal way. This must be wrapped in \makeatletter/\makeatother
% because \@ifclassloaded contains an @-catcode character.
% -----------------------------------------------------------------------------
\makeatletter
\@ifclassloaded{beamer}{%
  \usefonttheme{professionalfonts}
  \setbeamercolor{structure}{fg=primary-blue}
  \setbeamercolor{frametitle}{fg=primary-blue}
  \setbeamercolor{title}{fg=primary-blue}
  \setbeamercolor{subtitle}{fg=primary-gold}
  \setbeamercolor{author}{fg=primary-blue}
  \setbeamercolor{institute}{fg=neutral}
  \setbeamercolor{date}{fg=primary-gold}
  \setbeamercolor{itemize item}{fg=primary-blue}
  \setbeamercolor{itemize subitem}{fg=primary-gold}
  \setbeamercolor{alerted text}{fg=primary-gold}
  \setbeamercolor{block title}{fg=white, bg=primary-blue}
  \setbeamercolor{block body}{bg=light-bg}
  % Semantic block variants: block=definition/concept (blue), exampleblock=
  % worked example (green/positive), alertblock=key takeaway or caution (gold).
  % Keep to <=2 colored boxes per slide across ALL three variants (INV-7).
  \setbeamercolor{block title example}{fg=white, bg=positive}
  \setbeamercolor{block body example}{bg=light-bg}
  \setbeamercolor{block title alerted}{fg=white, bg=primary-gold}
  \setbeamercolor{block body alerted}{bg=light-bg}

  % Minimal footer: course code (left) + page X / N (right), no navigation clutter
  \setbeamertemplate{navigation symbols}{}
  \setbeamertemplate{footline}{%
    \color{neutral}\footnotesize\hspace{0.7em}\@coursecodetext\hfill
    \insertframenumber/\inserttotalframenumber\hspace{0.7em}\vspace{0.3em}}
}{}
\makeatother

% -----------------------------------------------------------------------------
% TikZ setup — libraries the snippet gallery uses
% -----------------------------------------------------------------------------
\usepackage{tikz}
\usetikzlibrary{arrows.meta, positioning, calc, decorations.pathreplacing,
                fit, shapes.geometric, backgrounds}

% Shared TikZ styles — snippets and hand-written diagrams can pick these up
% via `\begin{tikzpicture}[every node/.style={...}]` or by naming specific
% styles (e.g., `\node[dag-node] (X) at (0,0) {X};`).
\tikzset{
  % Node styles for causal DAGs and similar
  dag-node/.style = {
    draw=primary-blue, fill=light-bg, rounded corners=2pt,
    minimum width=1.2cm, minimum height=0.9cm, align=center, font=\small
  },
  decision-node/.style = {
    draw=primary-gold, fill=white, diamond, aspect=1.8,
    minimum width=1.8cm, minimum height=1.2cm, align=center, font=\small,
    inner sep=1pt
  },
  % Edge styles — observed vs counterfactual
  observed-edge/.style = {-{Latex[length=2.5mm]}, thick, draw=jet},
  counterfactual-edge/.style = {-{Latex[length=2.5mm]}, thick, draw=neutral, dashed},
  confound-edge/.style = {-{Latex[length=2.5mm]}, thick, draw=negative, dashed},
  % Data-point styles for plots
  observed-dot/.style = {fill=primary-blue, draw=primary-blue, circle, inner sep=1.2pt},
  counterfactual-dot/.style = {fill=white, draw=neutral, circle, inner sep=1.2pt}
}

% -----------------------------------------------------------------------------
% Convenience macros
% -----------------------------------------------------------------------------
\newcommand{\muted}[1]{\textcolor{neutral}{#1}}
\newcommand{\key}[1]{\textcolor{primary-gold}{\textbf{#1}}}
\newcommand{\good}[1]{\textcolor{positive}{#1}}
\newcommand{\bad}[1]{\textcolor{negative}{#1}}

% Standout frame macro: full-bleed dark background for section transitions.
% Beamer-only (uses \begin{frame}/\setbeamercolor/\usebeamerfont) -- do not
% call from an article-class file (e.g. Notes/<CODE>/*-notes.tex). \muted,
% \key, \good, \bad, and \coursecode above are plain xcolor/gdef and are
% safe to use from either class; verified when Notes/article-class support
% was added.
% Expand: \transitionslide{Your transition title}
\newcommand{\transitionslide}[1]{%
  \begingroup
  \setbeamercolor{background canvas}{bg=primary-blue}
  \begin{frame}[plain]
    \centering
    \vfill
    {\usebeamerfont{title}\color{white}\Large #1\par}
    \vfill
  \end{frame}
  \endgroup
}

% Section divider frame (Beamer-only), an alternative to \transitionslide with a
% darker two-line style: near-black (jet) background, a small white label line
% above a larger gold title, and a thin gold rule along the bottom edge.
% Expand: \sectiondivider{Section 2}{Pointers}
\newcommand{\sectiondivider}[2]{%
  \begingroup
  \setbeamercolor{background canvas}{bg=jet}
  \begin{frame}[plain]
    \vspace*{3.0cm}
    {\color{white}\ttfamily\large #1\par}
    \vspace{0.5cm}
    {\color{primary-gold}\LARGE #2\par}
    \begin{tikzpicture}[remember picture, overlay]
      \draw[primary-gold, line width=2.5pt]
        ($(current page.south west)+(0,0.1)$) -- ($(current page.south east)+(0,0.1)$);
    \end{tikzpicture}
  \end{frame}
  \endgroup
}

% -----------------------------------------------------------------------------
% Channel watermark — "Concepts That Click" (YouTube).
%
% Article-class files (CompetitiveExam Q/A sets, Notes, Assignments) get a
% light diagonal draft watermark on every page plus a centered footer naming
% the channel. Beamer decks get a subtle TikZ background watermark instead
% (the footline already carries the course code). Centralized here so every
% MCQ PDF is branded without per-file edits.
% -----------------------------------------------------------------------------
\makeatletter
\@ifclassloaded{article}{%
  \usepackage{draftwatermark}
  \SetWatermarkText{Concepts That Click}
  \SetWatermarkScale{1}
  \SetWatermarkFontSize{2cm}
  \SetWatermarkLightness{0.86}
  \SetWatermarkAngle{45}
  \usepackage{fancyhdr}
  \pagestyle{fancy}
  \fancyhf{}
  \fancyfoot[C]{\footnotesize\textcolor{neutral}{Concepts That Click~|~YouTube: \href{https://www.youtube.com/@concepts.thatclick}{@concepts.thatclick}}}
  \renewcommand{\headrulewidth}{0pt}
  \renewcommand{\footrulewidth}{0pt}
  \fancypagestyle{plain}{%
    \fancyhf{}%
    \fancyfoot[C]{\footnotesize\textcolor{neutral}{Concepts That Click~|~YouTube: \href{https://www.youtube.com/@concepts.thatclick}{@concepts.thatclick}}}%
    \renewcommand{\headrulewidth}{0pt}%
    \renewcommand{\footrulewidth}{0pt}%
  }
}{}
\@ifclassloaded{beamer}{%
  \setbeamertemplate{background}{%
    \begin{tikzpicture}[remember picture, overlay]
      \node[rotate=45, opacity=0.055, align=center]
        at (current page.center)
        {\Huge\textcolor{neutral}{Concepts That Click}};
    \end{tikzpicture}%
  }
}{}
\makeatother 
\coursecode{JKSSB}
\title{Lesson 65 Practice: Current Multimedia, GUI and Media Formats}
\author{JKSSB Computer Awareness}
\date{\today}
\begin{document}
\maketitle
\noindent\textbf{Provenance:} All 20 questions are original JKSSB-pattern
questions (\textit{[Original]}); none reproduces an official paper item.

\begin{enumerate}[label=\textbf{Q\arabic*.}, leftmargin=*]
\item \textit{[Original]} Expand GUI.
  \begin{enumerate}[label=\Alph*.]
    \item General User Interface
    \item Graphical User Interface
    \item Graphical Utility Interconnect
    \item Graphic User Interpreter
  \end{enumerate}
\item \textit{[Original]} Expand WYSIWYG.
  \begin{enumerate}[label=\Alph*.]
    \item What You See Is What You Get
    \item What You Save Is What You Get
    \item Where You See Is Where You Go
    \item What You Send Is What You Gain
  \end{enumerate}
\item \textit{[Original]} Expand MP3.
  \begin{enumerate}[label=\Alph*.]
    \item MPEG-3 Audio Layer I
    \item MPEG Audio Video Layer III
    \item Moving Picture Audio Layer III
    \item MPEG-1 Audio Layer III
  \end{enumerate}
\item \textit{[Original]} Expand PNG.
  \begin{enumerate}[label=\Alph*.]
    \item Portable Network Graphics
    \item Personal Network Graphics
    \item Portable Native Graphics
    \item Portable Network Group
  \end{enumerate}
\item \textit{[Original]} Expand JPEG.
  \begin{enumerate}[label=\Alph*.]
    \item Joint Picture Encoding Group
    \item Java Photographic Experts Group
    \item Joint Photographic Experts Group
    \item Joint Photographic Encoding Group
  \end{enumerate}
\item \textit{[Original]} Expand MP4.
  \begin{enumerate}[label=\Alph*.]
    \item MPEG-4 Portable 14
    \item MPEG-4 Picture 4
    \item MPEG-4 Player 4
    \item MPEG-4 Part 14
  \end{enumerate}
\item \textit{[Original]} In digitising an analogue signal, the process of
  measuring the signal at regular intervals is called:
  \begin{enumerate}[label=\Alph*.]
    \item quantization
    \item sampling
    \item synchronization
    \item compression
  \end{enumerate}
\item \textit{[Original]} The process of mapping each measured sample to one of
  a fixed set of discrete levels is called:
  \begin{enumerate}[label=\Alph*.]
    \item sampling
    \item synchronization
    \item quantization
    \item buffering
  \end{enumerate}
\item \textit{[Original]} Which of the following uses \textbf{lossless}
  compression?
  \begin{enumerate}[label=\Alph*.]
    \item JPEG
    \item MP3
    \item PNG
    \item MP4
  \end{enumerate}
\item \textit{[Original]} Which of the following correctly pairs an image
  format with its nature?
  \begin{enumerate}[label=\Alph*.]
    \item PNG --- lossy, does not support transparency
    \item JPEG --- lossless, best for photographs
    \item GIF --- lossless, limited to 256 colours
    \item BMP --- a small web format for photographs
  \end{enumerate}
\item \textit{[Original]} Which image format supports simple animation?
  \begin{enumerate}[label=\Alph*.]
    \item PNG
    \item GIF
    \item JPEG
    \item BMP
  \end{enumerate}
\item \textit{[Original]} Which format is most suitable for photographs?
  \begin{enumerate}[label=\Alph*.]
    \item PNG
    \item GIF
    \item JPEG
    \item TXT
  \end{enumerate}
\item \textit{[Original]} In multimedia, the alignment in time of separate
  streams such as audio and video is called:
  \begin{enumerate}[label=\Alph*.]
    \item compression
    \item sampling
    \item synchronization
    \item quantization
  \end{enumerate}
\item \textit{[Original]} The four classic elements of a GUI are remembered by
  the short form WIMP. Which set is correct?
  \begin{enumerate}[label=\Alph*.]
    \item Windows, Icons, Menus, Pointers
    \item Windows, Images, Mouse, Program
    \item Wires, Icons, Menus, Panels
    \item Windows, Icons, Modes, Printers
  \end{enumerate}
\item \textit{[Original]} A small window that asks the user for input or
  confirms an action, such as Save As or Print, is called a:
  \begin{enumerate}[label=\Alph*.]
    \item desktop
    \item dialogue box
    \item taskbar
    \item shortcut
  \end{enumerate}
\item \textit{[Original]} WYSIWYG means that:
  \begin{enumerate}[label=\Alph*.]
    \item the document is stored on the cloud.
    \item the screen shows the document exactly as it will be printed.
    \item the software corrects spelling automatically.
    \item the file is compressed before saving.
  \end{enumerate}
\item \textit{[Original]} A plotter is best suited for:
  \begin{enumerate}[label=\Alph*.]
    \item printing ordinary text letters.
    \item large engineering drawings and maps.
    \item scanning photographs into a computer.
    \item recording sound during a presentation.
  \end{enumerate}
\item \textit{[Original]} Which of the following is a \textbf{non-impact}
  printer?
  \begin{enumerate}[label=\Alph*.]
    \item Dot-matrix printer
    \item Daisy-wheel printer
    \item Line printer
    \item Laser printer
  \end{enumerate}
\item \textit{[Original]} Which of the following is \textbf{NOT} a lossless
  format?
  \begin{enumerate}[label=\Alph*.]
    \item PNG
    \item GIF
    \item JPEG
    \item ZIP
  \end{enumerate}
\item \textit{[Original]} Evaluate the statements:
  \begin{enumerate}[label=\Roman*.]
    \item MP3 is a lossy audio format.
    \item MP4 is a container format that can hold video and audio together.
    \item PNG is a lossless image format that supports transparency.
  \end{enumerate}
  \begin{enumerate}[label=\Alph*.]
    \item I and II only
    \item I and III only
    \item II and III only
    \item I, II, and III
  \end{enumerate}
\end{enumerate}
\end{document}
